\setcounter{section}{7}

\Needspace{5\baselineskip}
\hypertarget{ux7ebfux6027ux6ce8ux610fux529blinear-attention}{%
\section{线性注意力（Linear Attention）}\label{ux7ebfux6027ux6ce8ux610fux529blinear-attention}}

\Needspace{5\baselineskip}
\hypertarget{ux4e00ux57faux672cux539fux7406-1}{%
\subsection{一、基本原理}\label{ux4e00ux57faux672cux539fux7406-1}}

线性化注意力的核心思想是：通过改变计算顺序，避免显式地计算和存储这个 \(n \times n\) 的注意力矩阵 \(A\)。

它的目标是将计算复杂度从 \(O(n^2)\) 降低到 \(O(n)\)。

如何实现？答案是利用矩阵乘法的结合律。

\begin{itemize}
\tightlist
\item
  标准计算顺序：\((QK^T)V\)，先算 \(n \times n\) 矩阵，再与 \(V\)（\(n \times d\)）相乘。
\item
  线性化计算顺序：\(Q(K^T V)\)，先算 \(K^T V\)（得到一个 \(d \times d\) 的矩阵），再与 \(Q\)（\(n \times d\)）相乘。
\end{itemize}

\Needspace{5\baselineskip}
实现线性化并非无条件的，必须去掉 Softmax，这是因为：

\Needspace{5\baselineskip}
标准注意力（Standard Attention）为：

\[
\mathrm{Output} = \mathrm{Softmax}(QK^T)V
\]

请注意，这里不仅有乘法，还有一个非线性的 Softmax 包裹在中间。

\[
\mathrm{Softmax}(QK^T) \ne Q \cdot \mathrm{Softmax}(K^T)
\]

你不能把 Softmax 拆开，也不能把它挪到后面去先算 \(K^T V\)。因为 \(Q\) 和 \(K\) 的每一项都必须先交互、求和、归一化之后，才能和 \(V\) 相乘。这意味着你必须构建那个巨大的 \(n \times n\) 矩阵。

我们需要一个核函数来近似 Softmax 函数。具体数学推导如下。

\Needspace{5\baselineskip}
首先，我们把标准注意力写成一个按行归一化的形式。对于第 \(i\) 个查询向量 \(q_i\) 的输出 \(o_i\)：

\[
o_i =
\frac{
\sum_{j=1}^{n} \exp\left(\frac{q_i k_j^T}{\sqrt{d_k}}\right)v_j
}{
\sum_{j=1}^{n} \exp\left(\frac{q_i k_j^T}{\sqrt{d_k}}\right)
}
\]

\Needspace{5\baselineskip}
我们可以将其抽象为一个更通用的相似度函数 \(\mathrm{sim}(\cdot,\cdot)\)：

\[
o_i =
\frac{
\sum_{j=1}^{n} \mathrm{sim}(q_i,k_j)v_j
}{
\sum_{j=1}^{n} \mathrm{sim}(q_i,k_j)
}
\]

\Needspace{5\baselineskip}
接着引入核函数 \(\phi\)，用特征映射后的内积来表示相似度：

\[
\mathrm{sim}(q_i,k_j) = \phi(q_i)\cdot\phi(k_j)^T
\]

其中 \(\phi:\mathbb{R}^{d}\rightarrow\mathbb{R}^{d_{proj}}\) 是一个将原始向量映射到另一个特征空间的函数，它也称为核函数。

\Needspace{5\baselineskip}
现在，把这个核函数代入输出 \(o_i\) 的计算中：

\[
o_i =
\frac{
\sum_{j=1}^{n}[\phi(q_i)\cdot\phi(k_j)^T]v_j
}{
\sum_{j=1}^{n}[\phi(q_i)\cdot\phi(k_j)^T]
}
\]

\Needspace{5\baselineskip}
由于 \(\phi(q_i)\) 对于求和索引 \(j\) 是常数，我们可以将其提到求和号外面：

\[
o_i =
\frac{
\phi(q_i)\cdot\left[\sum_{j=1}^{n}\phi(k_j)^T \otimes v_j\right]
}{
\phi(q_i)\cdot\left[\sum_{j=1}^{n}\phi(k_j)^T\right]
}
\]

这里 \(\otimes\) 表示外积；在固定行、列向量方向后，也常写成矩阵乘法。

注：内积表示行向量乘列向量，结果为标量；外积表示列向量（\(m \times 1\)）乘行向量（\(1 \times n\)），结果为矩阵（\(m \times n\)）。

\Needspace{5\baselineskip}
让我们定义两个关键的累积状态矩阵：

\begin{enumerate}
\def\labelenumi{\arabic{enumi}.}
\tightlist
\item
  \(S\)：键和值的聚合状态。
\end{enumerate}

\[
S = \sum_{j=1}^{n}\phi(k_j)^T v_j \in \mathbb{R}^{d_{proj}\times d_v}
\]

注意：这里的 \(\phi(k_j)\) 是行向量，它转置后是列向量，与 \(v_j\)（行向量）做外积，得到一个矩阵。我们对所有 \(j\) 的这样的矩阵求和。

\begin{enumerate}
\def\labelenumi{\arabic{enumi}.}
\setcounter{enumi}{1}
\tightlist
\item
  \(Z\)：归一化因子的聚合状态。
\end{enumerate}

\[
Z = \sum_{j=1}^{n}\phi(k_j)^T \in \mathbb{R}^{d_{proj}\times 1}
\]

\Needspace{5\baselineskip}
现在，对于整个序列的输出 \(O\) 中的第 \(i\) 行 \(o_i\)，可以计算为：

\[
o_i = \frac{\phi(q_i)S}{\phi(q_i)Z}
\]

这里的分母 \(\phi(q_i)Z\) 是一个标量（归一化因子），因此输出向量的每个分量都除以同一个归一化值。

\begin{itemize}
\tightlist
\item
  计算 \(S\) 和 \(Z\)：遍历序列中的每个位置 \(j\)，求 \(\phi(k_j)^T v_j\) 和 \(\phi(k_j)^T\)，然后累加。这一步的复杂度是 \(O(n\cdot d_{proj}\cdot d_v)\)。
\item
  计算所有输出 \(o_i\)：对于每个位置 \(i\)，计算 \(\phi(q_i)S\) 和 \(\phi(q_i)Z\)。这一步的复杂度也是 \(O(n\cdot d_{proj}\cdot d_v)\)。
\end{itemize}

总复杂度是 \(O(nd_{proj}d_v)\)。如果我们固定 \(d_{proj}\) 和 \(d_v\)，它们是与序列长度 \(n\) 无关的常数，那么复杂度就是 \(O(n)\)，即线性复杂度。

注：\(K\) 和 \(V\) 本来分别为 \(n \times d_k\) 和 \(n \times d_v\) 维，\(K\) 转置后与 \(V\) 矩阵相乘，矩阵乘法计算复杂度为 \(n\cdot d_k\cdot d_v\)，也就是相对于序列长度的 \(O(n)\)。这里的 \(\phi\) 只改变特征维度，不改变对序列位置逐项累加的线性结构。

\Needspace{5\baselineskip}
\hypertarget{ux4e8cux5982ux4f55ux5bfbux627eux6838ux51fdux6570-phix}{%
\subsection{\texorpdfstring{二、如何寻找核函数 \(\phi(x)\)？}{二、如何寻找核函数 \textbackslash phi(x)？}}\label{ux4e8cux5982ux4f55ux5bfbux627eux6838ux51fdux6570-phix}}

\Needspace{5\baselineskip}
对于 Softmax 中的指数相似度，我们希望找到某种特征映射，使得：

\[
\mathrm{sim}(q,k) = \exp(q\cdot k^T) \approx \phi(q)\cdot\phi(k)^T
\]

在深度学习模型（如 Transformer 的原始注意力机制）中直接使用 Softmax 时，我们是在进行数值计算。对于单个数值 \(x\)，\(\exp(x)\) 是一个成熟的数学函数，计算机可以高效处理。但当 \(q\) 和 \(k\) 是两个向量时，若希望把 \(\exp(q\cdot k)\) 写成两个特征向量的内积，就会自然引出更高维的特征映射。

\Needspace{5\baselineskip}
从泰勒展开开始：

\[
\exp(q\cdot k)
= 1 + (q\cdot k) + \frac{(q\cdot k)^2}{2!}
+ \frac{(q\cdot k)^3}{3!}
+ \frac{(q\cdot k)^4}{4!}
+ \cdots
\]

\Needspace{5\baselineskip}
对于 \(m=1\)：

\[
(q\cdot k)^1 = \sum_{i=1}^{d}q_i k_i
\]

这可以看作是向量 \(q\) 和 \(k\) 在原始 \(d\) 维空间中的内积。

\Needspace{5\baselineskip}
对于 \(m=2\)：

\[
(q\cdot k)^2
= \left(\sum_{i=1}^{d}q_i k_i\right)^2
= \sum_{i=1}^{d}\sum_{j=1}^{d}q_i q_j k_i k_j
\]

\Needspace{5\baselineskip}
注意，这可以重写为两个新向量的内积：

\begin{itemize}
\tightlist
\item
  令 \(\phi_2(q)\) 是一个 \(d^2\) 维的向量，包含所有 \(q_i q_j\) 的组合。
\item
  令 \(\phi_2(k)\) 是一个 \(d^2\) 维的向量，包含所有 \(k_i k_j\) 的组合。
\end{itemize}

\Needspace{5\baselineskip}
那么：

\[
(q\cdot k)^2 = \phi_2(q)\cdot\phi_2(k)
\]

\Needspace{5\baselineskip}
现在，我们可以为每个幂次 \(m\) 定义一个特征映射：

\begin{itemize}
\tightlist
\item
  对于 \(m=0\)：\(\phi_0(q)=1\)，是 1 维。
\item
  对于 \(m=1\)：\(\phi_1(q)=q\)，是 \(d\) 维。
\item
  对于 \(m=2\)：\(\phi_2(q)\) 包含所有 \(q_i q_j\)，是 \(d^2\) 维。
\item
  对于 \(m=3\)：\(\phi_3(q)\) 包含所有 \(q_i q_j q_k\)，是 \(d^3\) 维。
\item
  依此类推。
\end{itemize}

\Needspace{5\baselineskip}
将这些不同阶数的特征组合起来，就可以构造一个无限维的超级特征映射 \(\Phi(q)\)：

\[
\Phi(q)=
\left(
1,\;
q,\;
\frac{q_i q_j}{\sqrt{2!}},\;
\frac{q_i q_j q_k}{\sqrt{3!}},\;
\frac{q_i q_j q_k q_l}{\sqrt{4!}},\;
\cdots
\right)
\]

这里每个分量都除以相应的 \(\sqrt{m!}\)，用于匹配泰勒展开中的系数。

\Needspace{5\baselineskip}
于是：

\[
\begin{aligned}
\Phi(q)\cdot\Phi(k)
&= 1\cdot 1
+ \sum_i q_i k_i
+ \sum_{i,j}\frac{q_i q_j}{\sqrt{2!}}\frac{k_i k_j}{\sqrt{2!}}
+ \sum_{i,j,k}\frac{q_i q_j q_k}{\sqrt{3!}}\frac{k_i k_j k_k}{\sqrt{3!}}
+ \cdots \\
&= 1 + \sum_i q_i k_i
+ \frac{1}{2!}\sum_{i,j}q_i q_j k_i k_j
+ \frac{1}{3!}\sum_{i,j,k}q_i q_j q_k k_i k_j k_k
+ \cdots
\end{aligned}
\]

\Needspace{5\baselineskip}
但注意：

\[
\sum_i q_i k_i = q\cdot k
\]

\[
\sum_{i,j}q_i q_j k_i k_j = (q\cdot k)^2
\]

\[
\sum_{i,j,k}q_i q_j q_k k_i k_j k_k = (q\cdot k)^3
\]

\Needspace{5\baselineskip}
因此：

\[
\Phi(q)\cdot\Phi(k)
= 1 + (q\cdot k) + \frac{(q\cdot k)^2}{2!}
+ \frac{(q\cdot k)^3}{3!}
+ \cdots
= \exp(q\cdot k)
\]

计算 \(\exp(q\cdot k)\) 这个看似简单的操作，在数学上等价于：将向量 \(q\) 和 \(k\) 映射到一个无限维的特征空间，在这个无限维空间中计算它们的内积。这个无限维空间包含了原始向量所有可能的多项式组合，包括一次项、二次项、三次项等。

线性注意力的出现本身就是为了解决长序列上 \(O(n^2)\) 的问题，但如果直接使用这个无限维特征映射，它自身又会带来趋于无限维的特征空间。在一个本身就可能趋于无限长的序列上，这是计算机无法处理的。

因此，我们必须对这个函数进行近似。常见近似方法如下。

\begin{enumerate}
\def\labelenumi{\arabic{enumi}.}
\item
  \(\phi(x)=\mathrm{elu}(x)+1\)。

  ELU 函数：对于 \(x>0\)，输出 \(x\)；对于 \(x\le 0\)，输出 \(\alpha(\exp(x)-1)\)，通常 \(\alpha=1\)。

  \(\mathrm{elu}(x)+1\)：对于 \(x>0\)，约为 \(x+1\)（线性增长）；对于 \(x\le 0\)，变为 \(\exp(x)\)（指数增长）。

  设计动机：这个函数在 \(x\le 0\) 时具有类似指数函数的行为，而在 \(x>0\) 时保持线性，整体上试图模拟指数函数的形状。但它是有限维的，输出维度与输入 \(x\) 相同。
\item
  \(\phi(x)=\mathrm{ReLU}(x)^2\)。

  这是一个简单的多项式函数。\(\mathrm{ReLU}(x)\) 将负值置零、正值保留，然后再平方。

  它对应一个二阶多项式核，计算简单，但近似能力相对较弱。
\item
  基于随机特征的方法。

  这是一种数学上更严谨的近似方法。它通过随机投影构造一个显式、有限维的特征映射 \(\phi(x)\)，使得 \(\phi(q)\cdot\phi(k)\) 近似于某个目标核函数。
\item
  学习得到的核函数。

  不手动设计 \(\phi\)，而是让模型自己学习。可以用一个小型神经网络（如一层 MLP）作为 \(\phi\)，在训练过程中通过梯度下降优化它，使其能更好地配合整个模型完成任务。
\end{enumerate}

但这些核函数都不能足够好地拟合 Softmax，这是线性注意力性能不如 Full Attention 的核心原因。

\Needspace{5\baselineskip}
\hypertarget{ux4e09ux5728ux81eaux56deux5f52ux4e2dux7684ux4f18ux52bf}{%
\subsection{三、在自回归中的优势}\label{ux4e09ux5728ux81eaux56deux5f52ux4e2dux7684ux4f18ux52bf}}

线性化注意力还有一个独特优势：它可以被轻松地改写为递归形式，非常适合自回归推理。

\Needspace{5\baselineskip}
令：

\[
S_t = \sum_{j=1}^{t}\phi(k_j)^T v_j
= S_{t-1}+\phi(k_t)^T v_t
\]

\[
Z_t = \sum_{j=1}^{t}\phi(k_j)^T
= Z_{t-1}+\phi(k_t)^T
\]

\Needspace{5\baselineskip}
在生成第 \(t\) 个 token 时：

\begin{enumerate}
\def\labelenumi{\arabic{enumi}.}
\tightlist
\item
  我们已经有之前所有 token 的累积状态 \(S_{t-1}\) 和 \(Z_{t-1}\)。
\item
\Needspace{5\baselineskip}
  我们计算当前 token 的 \(k_t\)、\(v_t\)，并更新状态：
\end{enumerate}

\[
S_t = S_{t-1}+\phi(k_t)^T v_t,\quad
Z_t = Z_{t-1}+\phi(k_t)^T
\]

\Needspace{4\baselineskip}
\begin{enumerate}
\def\labelenumi{\arabic{enumi}.}
\setcounter{enumi}{2}
\tightlist
\item
\Needspace{5\baselineskip}
  然后使用当前查询 \(q_t\) 计算输出：
\end{enumerate}

\[
o_t = \frac{\phi(q_t)S_t}{\phi(q_t)Z_t}
\]

这就像一个 RNN。在生成第 \(t\) 个 token 时，我们已经有了前 \(t-1\) 个 token 的完整总结 \(S_{t-1}\) 和 \(Z_{t-1}\)。当新的第 \(t\) 个 token 产生时，我们不需要回溯整个历史，只需要将当前 token 的贡献 \(\phi(k_t)^T v_t\) 加到 \(S_{t-1}\) 上，将 \(\phi(k_t)^T\) 加到 \(Z_{t-1}\) 上，就得到了新的状态 \(S_t\) 和 \(Z_t\)。然后用当前查询 \(q_t\) 与最新状态交互，得到当前时刻的输出。

因此，我们不需要存储整个键值的历史，只需要维护一个固定大小的状态 \(S\) 和 \(Z\)。每一步的计算成本是常数 \(O(1)\)（相对于序列长度），极大地提高了长文本生成的效率。

\FloatBarrier
\Needspace{5\baselineskip}
\hypertarget{ux53c2ux8003ux6587ux732e-98}{%
\subsection{参考文献}\label{ux53c2ux8003ux6587ux732e-98}}

\vspace{0.25em}
\begin{itemize}
\tightlist
\item
  Katharopoulos, A., Vyas, A., Pappas, N., \& Fleuret, F. (2020). \href{https://arxiv.org/abs/2006.16236}{Transformers are RNNs: Fast Autoregressive Transformers with Linear Attention}. ICML 2020.
\item
  Shen, Z., Zhang, M., Zhao, H., Yi, S., \& Li, H. (2021). \href{https://arxiv.org/abs/1812.01243}{Efficient Attention: Attention with Linear Complexities}. WACV 2021.
\end{itemize}

\clearpage
